\subsection{Introducción del tema}

ROS (\textit{Robot Operating System}) es un \textit{framework} de código
abierto para el desarrollo de software robótico. A pesar de su nombre, no
constituye un sistema operativo en sentido estricto: se ejecuta sobre un
sistema operativo anfitrión, principalmente Linux, y provee una capa de
\textit{middleware} junto con un conjunto de bibliotecas y herramientas. Su
función central es permitir que los distintos componentes de un robot se
desarrollen como procesos independientes y se comuniquen entre sí mediante
interfaces estandarizadas.

El proyecto se originó en 2007 en el Laboratorio de Inteligencia Artificial de
la Universidad de Stanford y continuó su desarrollo en Willow Garage.
Actualmente es mantenido por la Open Source Robotics Foundation (OSRF) junto a
una amplia comunidad internacional.

Las principales razones de su adopción son:
\begin{itemize}
    \item \textbf{Modularidad:} cada función del robot se implementa como un
    componente independiente, lo que facilita el desarrollo en paralelo y el
    mantenimiento.
    \item \textbf{Reutilización:} existe un amplio repositorio de paquetes para
    navegación, percepción, manipulación y controladores de hardware.
    \item \textbf{Abstracción de hardware:} los componentes se comunican a
    través de interfaces definidas, independientemente del fabricante del
    sensor o actuador.
    \item \textbf{Herramientas:} incluye utilidades de visualización, registro,
    simulación y depuración.
    \item \textbf{Licenciamiento abierto:} se distribuye bajo licencias
    permisivas (BSD y Apache~2.0), lo que habilita su uso académico e
    industrial.
\end{itemize}

ROS~1 fue concebido como una plataforma de investigación, orientada a un único
robot operando en una red local confiable. A medida que su uso se extendió a
entornos industriales, sistemas multi-robot y redes inalámbricas, sus
limitaciones de diseño se hicieron evidentes. En respuesta, se desarrolló
ROS~2, rediseñado desde su base sobre estándares industriales de
comunicaciones.

La diferencia fundamental entre ambas versiones reside en su arquitectura de
comunicaciones. ROS~1 emplea un esquema híbrido en el que un nodo maestro
central (\textit{roscore}) gestiona el descubrimiento de los participantes,
mientras que el intercambio de datos se realiza mediante protocolos propios
(TCPROS y UDPROS). Este diseño introduce un punto único de falla y no ofrece
mecanismos para ajustar el comportamiento de la comunicación según las
condiciones de la red.

ROS~2, en cambio, adopta una arquitectura descentralizada basada en el
estándar DDS (\textit{Data Distribution Service}) y su protocolo de
interoperabilidad RTPS. Esta elección elimina la dependencia de un nodo
central y habilita la configuración de políticas de calidad de servicio (QoS),
el soporte para sistemas operativos de tiempo real y la integración de
microcontroladores mediante micro-ROS. Asimismo, incorpora mecanismos de
seguridad basados en la especificación DDS-Security, ausentes en ROS~1.

Con la finalización del soporte de ROS~1 en 2025, ROS~2 se ha consolidado como
la plataforma vigente para el desarrollo de nuevos sistemas robóticos.

Los aspectos técnicos de cada arquitectura se desarrollan en detalle en el
Capítulo~3.

\subsection{Descripción del contexto del tema}

ROS se utiliza como capa de integración en sistemas robóticos de diversa
complejidad, desde plataformas educativas y de investigación hasta vehículos
autónomos y robots industriales. En estos sistemas, los componentes de
software suelen distribuirse entre varios equipos: computadoras embarcadas en
el robot, microcontroladores asociados a sensores y actuadores, y estaciones de
supervisión remotas. ROS actúa como el \textit{middleware} que permite que
todos estos componentes intercambien información de forma transparente, sin
importar en qué equipo se ejecuten.

El software de ROS se publica en \textbf{distribuciones}: conjuntos versionados
de paquetes que se prueban y se mantienen en conjunto. Desde 2017, ROS~2 lanza
una nueva distribución cada año, el 23 de mayo. Las distribuciones de los años
pares son de soporte extendido (LTS, \textit{Long Term Support}) y reciben
actualizaciones durante cinco años, mientras que las de los años impares tienen
un ciclo de soporte de aproximadamente un año y medio. En paralelo existe
Rolling, una distribución de desarrollo continuo en la que se incorporan los
cambios antes de su inclusión en una versión estable.

Cada distribución está vinculada a una versión LTS específica de Ubuntu, que
constituye su plataforma principal de soporte (\textit{Tier~1}), según lo
establecido en la REP~2000. Este vínculo se debe a que los paquetes binarios de
ROS se compilan y se prueban contra las bibliotecas del sistema de esa versión
de Ubuntu. En consecuencia, la elección de la distribución de ROS determina la
versión del sistema operativo a utilizar. Las distribuciones LTS más recientes
y su correspondencia son:
\begin{itemize}
    \item \textbf{ROS~1 Noetic Ninjemys} (2020): Ubuntu 20.04. Última
    distribución de ROS~1, con fin de soporte en mayo de 2025.
    \item \textbf{ROS~2 Humble Hawksbill} (2022): Ubuntu 22.04.
    \item \textbf{ROS~2 Jazzy Jalisco} (2024): Ubuntu 24.04.
    \item \textbf{ROS~2 Lyrical Luth} (2026): Ubuntu 26.04.
\end{itemize}

Además de Ubuntu, las distribuciones de ROS~2 ofrecen distintos niveles de
soporte para otras plataformas, como Windows y Red Hat Enterprise Linux.

\subsection{Mercado y Segmentación del tema}

Como es el mercado de usabilidad y la segmentación del mismo. Como se usa
en el mundo y argentina

\subsection{Legislación del tema}
Si existen leyes o normas sobre el tema, el contexto mundial o de argentina